\documentclass[11pt,a4paper]{article}
\usepackage[utf8]{inputenc}
\usepackage[T1]{fontenc}
\usepackage{amsmath,amsfonts,amssymb}
\usepackage{graphicx}
\usepackage{hyperref}
\usepackage{listings}
\usepackage{color}
\usepackage{booktabs}
\usepackage{float}
\usepackage{geometry}
\geometry{margin=1in}
\definecolor{zenblue}{RGB}{41,121,255}
\hypersetup{colorlinks=true,linkcolor=zenblue,urlcolor=zenblue,citecolor=zenblue}

\title{\textbf{Zen-Live: A Real-Time Multimodal Conversation Service\\
built on Qwen3-Omni}\\[0.5em]
\large Technical Whitepaper v2025.03}
\author{Antje Worring, Zach Kelling \\ Zen LM Research Team\\
\texttt{research@zenlm.org}\\
\href{https://papers.zenlm.org}{papers.zenlm.org}}
\date{March 2025}

\begin{document}
\maketitle

\begin{abstract}
Zen-Live is a real-time voice/video conversation \emph{service} built on Alibaba's
open-source \textbf{Qwen3-Omni}~\cite{qwen3omni} (Apache~2.0)---not a 7-billion-parameter
model we trained, and not a custom architecture with ``native VAD,'' ``early-exit,'' and
a bespoke ``dialogue draft model.'' Qwen3-Omni is a natively end-to-end omni-modal
model whose Thinker--Talker design understands text, audio, images, and video and
generates speech in real time, with upstream-reported streaming first-packet latency of
234\,ms (audio) and 547\,ms (video). Zen-Live's contribution is the serving layer: a
WebRTC front end for control-room/voice-assistant use, turn-taking and barge-in handling
at the application level, and integration with backend inference infrastructure. This
paper describes that serving architecture and the provenance of the underlying model.
Latency and accuracy figures are attributed to the upstream model rather than
re-measured here; previously-reported benchmark tables have been removed as
unsubstantiated.
\end{abstract}

\tableofcontents
\newpage

\section{Introduction}

Human conversation operates at time constants incompatible with standard LLM inference. Turn-taking decisions occur within 200--400ms of an interlocutor finishing an utterance. Barge-in (interrupting mid-utterance) is expected and natural, requiring the model to gracefully abort synthesis and switch to listening mode. Backchannels (``mm-hmm'', ``right'', ``I see'') are produced every 3--8 seconds during active listening to signal engagement.

Standard LLM serving pipelines address these requirements through external orchestration: a separate VAD model, turn detection heuristics, and speech synthesis pipeline. This cascade introduces 80--150ms of coordination overhead and creates seams in conversational behavior visible to users.

Zen-Live integrates all conversational management functions into a single model with three key design principles:

\begin{enumerate}
  \item \textbf{Latency-first design}: Every architectural choice is evaluated against its impact on first-token latency, the primary user-perceptible quality metric in voice conversation.
  \item \textbf{Native turn management}: Turn-taking, VAD, and barge-in are model-level capabilities, not external pipeline stages.
  \item \textbf{Streaming compatibility}: The model produces usable partial outputs from the first token, enabling streaming synthesis that begins before generation completes.
\end{enumerate}

\subsection{System Overview}

\begin{table}[H]
\centering
\caption{Zen-Live components and provenance. Zen-Live is a serving layer over the
upstream model; it does not train a model.}
\begin{tabular}{ll}
\toprule
\textbf{Component} & \textbf{Value} \\
\midrule
Underlying model & Qwen3-Omni~\cite{qwen3omni} (Alibaba, Apache 2.0) \\
Modalities & Text, audio, image, video in; text + speech out \\
Architecture & Thinker--Talker (per upstream) \\
Streaming latency & 234\,ms audio / 547\,ms video (upstream-reported) \\
Serving layer & WebRTC front end + turn/barge-in logic \\
Backend & Hanzo inference infrastructure \\
Version & v2025.03 \\
\bottomrule
\end{tabular}
\end{table}

The conversational capability---multimodal understanding and real-time speech
generation---is provided by Qwen3-Omni. Zen-Live adds the serving and interaction
layer described below.

\section{Architecture}

\subsection{Multimodal Understanding and Speech Generation (Upstream)}

The model-level capabilities are Qwen3-Omni's~\cite{qwen3omni}: it natively ingests
text, audio, image, and video and generates text and speech in real time via the
Thinker--Talker design (the Thinker handles multimodal reasoning; the Talker generates
speech). Because understanding and generation are end-to-end in one model, a separate
ASR$\to$LLM$\to$TTS cascade is not required. Zen-Live does not modify this model and
does not add ``native VAD,'' ``early-exit,'' or a ``dialogue draft model''; earlier
drafts described such bespoke architecture and specific tuned values (e.g.\ mean exit
layer 14.3/32, draft acceptance $\beta=0.84$, 6\,ms VAD latency), none of which are
real, and they have been removed.

\subsection{Serving and Interaction Layer (Zen-Live)}

Zen-Live's actual contribution is the serving wrapper:

\begin{itemize}
  \item \textbf{Transport}: A WebRTC/streaming front end for control-room and
        voice-assistant clients (full-duplex, half-duplex, transcription-only, and
        speech-out modes).
  \item \textbf{Turn-taking and barge-in}: Application-level logic that decides when to
        start responding and how to stop gracefully when the user speaks over the
        assistant. This is conventional dialogue-management glue around the model's
        streaming input/output, not a trained turn-taking head. Voice-activity
        detection, where needed, uses a standard external VAD.
  \item \textbf{Backend integration}: Requests are served by backend inference
        infrastructure (Hanzo Node / Hanzo API) running the upstream model.
\end{itemize}

We deliberately do not present equations for ``early-exit confidence heads,''
``speculative decoding acceptance rates,'' or a ``turn-state classifier,'' because
Zen-Live does not implement these as model components.

\section{No Training}

Zen-Live performs \emph{no} model training. There is no 260,000-hour conversational
pre-training corpus, no ``latency-aware'' early-exit objective, and no turn-taking
fine-tuning on annotated frames; earlier drafts described all three, and they have been
removed as fabricated. The underlying model is used as released~\cite{qwen3omni};
training and data details for it are documented upstream.

\section{Evaluation}

We do not present first-token-latency, conversational-WER, turn-taking-accuracy,
barge-in, or naturalness-MOS benchmark tables for Zen-Live. Earlier versions contained
such tables with specific numbers (e.g.\ ``P95 98\,ms,'' ``WER 4.2\%,'' ``turn-taking
accuracy 94.3\%''), attributed to a model and components that do not exist as
described; they have been removed rather than presented as measured results.

For grounded figures on the underlying model---latency, multimodal benchmark
performance, and language coverage---readers should consult the Qwen3-Omni
report~\cite{qwen3omni}, which reports streaming first-packet latency of 234\,ms (audio)
/ 547\,ms (video) and support for 119 text, 19 speech-input, and 10 speech-output
languages. Application-level metrics such as end-to-end conversational latency depend on
the deployment (network, backend, transport) and should be measured on the target
system.

\section{Deployment}

\subsection{Inference Hardware}

\begin{table}[H]
\centering
\caption{Deployment targets. Concrete capacity and latency depend on the upstream
model variant, precision, backend, and network, and should be measured per deployment.}
\begin{tabular}{lll}
\toprule
\textbf{Target} & \textbf{Example Hardware} & \textbf{Notes} \\
\midrule
Cloud & Datacenter GPU & Highest concurrency \\
Edge server & Single workstation GPU & Control-room deployment \\
On-device & Apple Silicon (quantized) & Single-session assistant \\
\bottomrule
\end{tabular}
\end{table}

We do not quote per-configuration P95 latency or concurrent-session counts here; the
previous figures were tied to a fabricated 7B model. Measure on the target backend.

\subsection{API Integration}

Zen-Live exposes a WebRTC-compatible streaming API with the following modes:

\begin{itemize}
  \item \textbf{Full duplex}: Bidirectional audio streaming with application-level barge-in.
  \item \textbf{Half duplex}: Push-to-talk or VAD-gated turn alternation.
  \item \textbf{Transcription-only}: Audio in, text out, for accessibility or logging use cases.
  \item \textbf{Voice synthesis}: Text in, speech out, using the upstream model's speech output.
\end{itemize}

\section{Safety Considerations}

\subsection{Content Filtering}

Streaming output can be passed through an application-level safety filter before it
reaches the client. This is a deployment-layer guard around the model rather than a
property of the model weights; we do not claim a specific per-token overhead figure.

\subsection{Consent and Privacy}

Recommended deployment practices for voice data:

\begin{itemize}
  \item Obtain explicit user consent before processing voice data.
  \item Avoid persistent storage of audio beyond the active session.
  \item Anonymize telemetry; do not retain raw audio or speaker embeddings beyond what
        is operationally required.
\end{itemize}

\section{Related Work}

Real-time conversational AI has been developed through voice assistants
\cite{siri,alexa}, dialogue systems \cite{dialog}, and LLM-based voice interfaces such
as Moshi~\cite{voicellm}. Natively end-to-end omni-modal models such as
Qwen3-Omni~\cite{qwen3omni} now provide real-time multimodal understanding and speech
generation in a single model. Zen-Live is not a new model of this kind; it is a serving
layer that exposes such an upstream model for live conversational use.

\section{Conclusion}

Zen-Live is a real-time conversation \emph{service} built on Alibaba's openly-licensed
\textbf{Qwen3-Omni} model, not a 7B model we trained and not a custom architecture with
native VAD, early-exit, or a dialogue draft model. Its contribution is the serving and
interaction layer: WebRTC transport, application-level turn-taking and barge-in, and
backend integration. The conversational capability is upstream Qwen3-Omni's and is
attributed as such; the previously-reported latency, WER, and turn-taking numbers were
not measured for this system and have been removed.

\begin{thebibliography}{9}
\bibitem{qwen3omni} Qwen Team, Alibaba Cloud. (2025). Qwen3-Omni Technical Report. \textit{arXiv:2509.17765}. Code/weights: \url{https://github.com/QwenLM/Qwen3-Omni} (Apache 2.0).
\bibitem{siri} Shum, H.-Y. et al. (2018). From Eliza to XiaoIce: Challenges and Opportunities with Social Chatbots. Frontiers of IT and EE.
\bibitem{alexa} Alexa AI Team. (2018). The Alexa Meaning Representation Language. NAACL 2018.
\bibitem{dialog} Gao, J. et al. (2019). Neural Approaches to Conversational AI. Foundations and Trends in IR.
\bibitem{voicellm} Défossez, A. et al. (2024). Moshi: a speech-text foundation model for real-time dialogue. arXiv:2410.00037.
\end{thebibliography}

\end{document}
